% CARE 2.0 supplementary-information source for the Springer Nature template.
% Numerical values are copied from frozen result artifacts in the CARE2.0 repo.
\documentclass[pdflatex,sn-nature]{sn-jnl}

\usepackage{graphicx}
\usepackage{booktabs}
\usepackage{tabularx}
\usepackage{array}
\usepackage{amsmath}
\usepackage{amssymb}
\usepackage{xcolor}
\usepackage{url}
\usepackage{listings}
\usepackage{float}

\newcommand{\care}{CARE 2.0}
\newcommand{\dauc}{\Delta\mathrm{AUC}}
\newcommand{\dfinal}{\Delta\mathrm{Final}}
\newcommand{\source}{\mathcal{D}_{s}}
\newcommand{\target}{\mathcal{D}_{t}}
\newcommand{\code}[1]{\texttt{#1}}

\lstdefinestyle{carecode}{
  basicstyle=\footnotesize\ttfamily,
  breaklines=true,
  columns=fullflexible,
  frame=single,
  framerule=0.3pt,
  showstringspaces=false
}

\title[CARE 2.0 Supplementary Information]{Supplementary Information for CARE 2.0: Evidence-Bounded LLM Skill Transfer for Sequential Scientific Experiments}

\author*[1]{\fnm{Anonymous} \sur{Authors}}
\affil*[1]{\orgdiv{Anonymous submission}}

\begin{document}
\maketitle

\begin{abstract}
\mbox{T}his Supplementary Information reports the experimental protocol, full result
tables, matched baselines, ablations, decision traces, and claim boundaries for
\care. The system uses a large language model (LLM) in two distinct roles. An
offline semantic-skill compiler converts source-task evidence and a public
target schema into executable features, priors, and acquisition schedules. A
separate online proposer--critic controller may update a falsifiable transfer
hypothesis and select an experiment from a bounded candidate menu. We report
these roles separately because a model call, a plausible explanation, and a
causal improvement in the executed experiment are not equivalent. Frozen
held-out replay provides positive examples in molecular-property, reaction-HTE,
and materials-property tasks. Matched mechanism controls show that some of the
largest source-conditioned gains are explained by the initial experimental
design rather than continuous online transfer. A route-disjoint audit and a
six-route live pilot further show that online LLM participation does not yet
provide a stable general advantage. A newly frozen protein-engineering task
family produces strong negative transfer for all three classical transfer
baselines; one live LLM trajectory identifies the mismatch and hands control
back to target-only GP, but its initial design remains worse than a fixed
design. A retrospective gate calibrated only on completed source backbones
would reject all three transfer policies before reading P06241 outcomes and
avoid the 8.133 AUC loss of the least harmful transfer policy. Because this gate
was designed after the external result was known, it is a mechanism audit, not
prospective safety confirmation. We therefore preregistered the identical
source-only eligibility rule, target-only fallback, code, and seed schedule on
a second external protein family before downloading its official data. On the
untouched Rhodopsin test partition, the ungated source-only choice lost 1.548
AUC [95\% CI $-1.996,-1.099$]; the frozen gate rejected transfer and avoided
that loss. This is prospective negative-transfer prevention on one external
family. We then froze an outcome-blind LLM hypothesis on the official FLIP2
Imine Reductase two-to-many split before downloading the data. The LLM-generated
additive-mutation skill passed train-to-validation calibration and improved AUC
on the untouched test partition by 0.389 [95\% paired-bootstrap CI
$0.274,0.503$] relative to target-only GP-UCB. A predeclared fixed skill prior
improved more ($+0.864$), so this is prospective positive-transfer evidence for
one LLM-generated skill, not evidence that the LLM is the strongest optimizer.
A post-hoc audit against 99 deterministic source-outcome permutations found
that the true binding exceeded the permutation mean by 0.340 AUC, but the
assignment-level randomization test did not reject the null at 0.05
($p=0.07$). Thus the prospective efficacy result stands, whereas the proposed
outcome-binding mechanism remains suggestive and non-unique. Neither result is
wet-lab validation. We therefore froze efficacy and source-assignment endpoints
together on a new FLIP2 TrpB family before downloading its data. The
LLM-selected additive skill improved official-test AUC by 0.949 [95\% CI
$0.864,1.034$], and the true source binding exceeded the mean of 99 outcome
permutations by 0.890 AUC [95\% bootstrap CI $0.810,0.970$], with no null
assignment as large as the true effect ($p=0.01$). The frozen gate nevertheless
abstained because its single-mutant calibration interval crossed zero; exact
mutation-token coverage was 0\% on calibration but 99.97\% at deployment.
Thus TrpB supplies independent prospective efficacy and mechanism confirmation
for one LLM-selected skill while exposing a power failure in the present
router. It does not establish universal cross-domain transfer or wet-lab
discovery. All negative and null results are retained.
\end{abstract}

\begin{appendices}

\clearpage
\section{Supplementary evidence map}
\label{sec:evidence_map}

The appendix follows the same evidence-first organization used for CARE 1.0:
headline claims are linked to complete result tables, mechanism controls,
representative traces, and executable reproduction artifacts. Table
\ref{tab:evidence_tiers} prevents results from different experimental stages
from being pooled into one headline number.

\begin{table}[t]
\centering
\caption{Evidence tiers used throughout the supplement. ``Frozen'' means that
the skill or controller and the comparison rule were fixed before the reported
evaluation seeds.}
\label{tab:evidence_tiers}
\scriptsize
\setlength{\tabcolsep}{3pt}
\renewcommand{\arraystretch}{0.92}
\begin{tabularx}{\textwidth}{p{2.5cm}p{3.6cm}X p{3.2cm}}
\toprule
Tier & Main question & Evidence & Permitted interpretation \\
\midrule
Frozen semantic skill & Can an LLM-compiled representation improve a strong target-only optimizer? & 300--500 paired held-out seeds per target after development and calibration & Positive task-specific evidence for compiled skills \\
Mechanism audit & Which component produces the observed gain? & Same-seed schedule-only, no-prior, warm-start-only, and data-only controls & Retrospective component attribution \\
Hidden-label non-interference & Can unrevealed target labels affect a decision state? & Exact shadow-label reconstruction of 170 archived states across 17 routes & Implementation-boundary evidence, not efficacy evidence \\
Online controller audit & Does the LLM change decisions after a matched initial design? & Same-start target GP, RGPE, multitask GP, route inventory, and controller handoff replays & Development and robustness evidence, not external confirmation \\
Probability calibration audit & Does the LLM know when its executed candidate will improve the current best? & All 170 archived decisions, scored against the same-candidate GP probability & Retrospective uncertainty audit; no prospective gate claim \\
\bottomrule
\end{tabularx}
\end{table}

\begin{table}[t]
\centering
\caption{Evidence tiers used throughout the supplement (continued).}
\label{tab:evidence_tiers_continued}
\scriptsize
\setlength{\tabcolsep}{3pt}
\renewcommand{\arraystretch}{0.92}
\begin{tabularx}{\textwidth}{p{2.5cm}p{3.6cm}X p{3.2cm}}
\toprule
Tier & Main question & Evidence & Permitted interpretation \\
\midrule
Live bounded-authority pilot & Does one real LLM decision add value on disjoint routes? & Six successful Opus trajectories plus one preserved infrastructure failure & Operational pilot only; no confirmatory inference \\
External task-family stress test & Does a frozen controller generalize to a task family absent from development? & Official FLIP2 wild-type holdout; 100 paired classical-transfer seeds and one live Opus trajectory & Negative-transfer boundary and abstention case; not positive external confirmation \\
Source-only family gate & Can completed source tasks reject unsafe transfer without target outcomes? & Two 50-seed source-only pseudo-target routes, frozen lower-confidence-bound eligibility rule, and outcome-blind P06241 deployment replay & Retrospective mechanism audit; the rule was designed after the external result and is not prospective confirmation \\
Prospective family gate & Does the source-only rule prevent harm when frozen before a second external family is downloaded? & Preregistered FLIP2 Rhodopsin by-wild-type split; two 50-seed calibration routes and 100 untouched test seeds & Prospective negative-transfer prevention on one external family; not positive-transfer efficacy \\
Prospective positive-transfer skill & Can an outcome-blind LLM generate a useful skill before a new external dataset is downloaded? & Preregistered FLIP2 IRED two-to-many split; 100 train-to-validation calibration seeds, 100 untouched official-test seeds, and all paired baselines & Positive finite-pool efficacy for one LLM-generated additive skill; a fixed skill prior remains stronger \\
Prospective efficacy and attribution replication & Does an independently frozen LLM-selected skill improve a new family, and does that gain require the measured source binding? & Preregistered FLIP2 TrpB one-to-many split; 100 deployment seeds, 99 predeclared outcome permutations, and a sequence-only finite-pool rule & Positive skill efficacy and rejection of the source-assignment null; the frozen router abstained and did not realize the gain \\
Frozen repeated protocol & Are effects stable across stochastic model calls? & Declared 30 trajectories per route, hash-locked analysis, no optional stopping & Pending until all declared trajectories complete \\
\bottomrule
\end{tabularx}
\end{table}

The strongest current claim is deliberately narrower than ``a general AI
scientist improves every domain.'' The supported claim is that an LLM can
compile scientific descriptions into explicit skills that are executable,
calibrated, rejectable, and auditable; some of these skills improve finite-pool
search on held-out real datasets, and one preregistered gate prevented a harmful
transfer on a second external family. A separately preregistered IRED experiment
also supplies positive external efficacy for one outcome-blind LLM-generated
skill. A second preregistered TrpB family independently confirms positive skill
efficacy and source-outcome attribution. The present evidence does not establish
universal cross-domain transfer, elimination of negative transfer, or a new
physical discovery; the TrpB result also shows that the current gate can abstain
when a useful transfer is available.

\section{Problem setting and evaluation protocol}
\label{sec:protocol}

\subsection{Budgeted source-to-target transfer}

Let $\source=\{(x_i,y_i)\}$ denote a completed source experiment and let
$\target=\{x_j\}$ denote a finite target candidate pool. At target round $t$,
the system observes outcomes only for the selected set $O_t$ and chooses one
unmeasured candidate $x_t$. The target budget $B$ is fixed across CARE and all
comparators. The primary trajectory metric is the mean best-so-far value,

\begin{equation}
  \mathrm{AUC}(\pi)=\frac{1}{B}\sum_{t=1}^{B}
  \max_{x\in O_t} y(x),
\end{equation}

where larger values indicate that a policy finds strong candidates earlier.
We also report the final best value and whether a globally top-10 candidate was
observed. For transformed materials objectives, optimization is performed on
the declared dimensionless replay score; the raw property, physical unit,
transformation, and direction are retained in the outcome contract.

\subsection{LLM roles}
\label{sec:llm_roles}

The offline compiler receives source evidence and the public target schema. It
returns a small library of semantic features, optional coefficient directions,
and an acquisition schedule. A strict compiler rejects unknown fields, unknown
categorical values, invalid operators, and malformed output. Calibration uses
development seeds only. After selection, the skill is frozen and held-out
replay makes no model call.

The online controller is evaluated separately. At every permitted round it
receives: (i) the current transfer hypothesis, (ii) revealed target outcomes,
(iii) target-only GP predictions, and (iv) a bounded menu containing target-GP,
source-prior, consensus, and coverage candidates. A proposer recommends one
candidate and a critic may accept, revise, or return to GP rank one. The trace
stores both responses, the executed candidate, the hidden value revealed after
execution, prediction error, token use, and fallback state.

\subsection{Leakage boundary}

During skill generation, the LLM may see source summaries and public target
field names, value vocabularies, and counts. It does not see unrevealed target
outcomes, held-out candidate identifiers, answer keys, or comparator results.
At online round $t$, only outcomes revealed in rounds $<t$ are included. The
target-only comparator starts from the same target observations when isolating
the online-controller increment. Calibration, held-out replay, and online
trajectory evidence are stored in separate directories with immutable hashes.

\subsection{Empirical hidden-label non-interference}
\label{sec:noninterference}

We tested the information boundary through the production decision-state code
rather than relying only on a prompt-field blacklist. For every one of 170
archived online request states spanning 17 source--target routes, we preserved
the outcomes already revealed before that request, every public candidate
attribute, source evidence, and all protocol settings. We then permuted every
remaining hidden target label and rebuilt the source prior, candidate menu,
menu diagnostics, and complete LLM prompt. A state passed only if all four
objects were exactly invariant under canonical JSON serialization and the
archived prompt still matched its recorded SHA-256.

All 170 states passed, covering 40 molecular-property states, 80
materials-property states, and 50 reaction-optimization states. Across these
state-level interventions, 2,725,689 unrevealed label positions were permuted.
Figure~\ref{fig:hidden_label_noninterference} reports every route and round.

\begin{figure*}[t]
\centering
\includegraphics[width=0.98\textwidth]{hidden_target_noninterference.png}
\caption{Hidden-target non-interference audit generated from the archived
decision states. Panel A shows the number of unrevealed labels permuted at each
route and round; Panel B aggregates intervention coverage by route. Every state
preserved the source prior, executable candidate menu, diagnostics, and full
current-code LLM prompt exactly. Historical prompt hashes were verified
separately. All archived states predate the added public outcome-semantics
field; 110 states reproduce after removing only that field, whereas 60 states
from six earlier routes also predate the current eligibility and safety fields.
Exact historical byte reconstruction is therefore reported as code-version
drift rather than used as the primary endpoint.}
\label{fig:hidden_label_noninterference}
\end{figure*}

This audit supports the implementation-level claim that unrevealed target
labels do not affect the current decision-state construction on the tested
routes. It does not establish that an LLM decision is scientifically correct,
that transfer improves the target experiment, or that the same property holds
for an untested future implementation.

\subsection{Baselines and why they are included}

\begin{table}[t]
\centering
\caption{Comparator portfolio. Each baseline tests a different alternative
explanation for an apparent CARE gain.}
\label{tab:baselines}
\begin{tabularx}{\textwidth}{p{3.2cm}X X}
\toprule
Comparator & What it does & Why it is necessary \\
\midrule
Random / no-transfer & Samples without source evidence or follows the frozen target-only initialization & Establishes a weak lower bound and detects gains due only to added computation \\
Target GP-UCB & Fits a GP to revealed target outcomes and maximizes posterior mean plus exploration bonus & Strong matched target-only sequential optimizer \\
Mixed-kernel GP-EI & Uses numeric and categorical kernels with expected improvement & Strong target-only baseline for heterogeneous reaction and materials spaces \\
Target acquisition portfolio & Selects the strongest frozen target-only acquisition policy & Conservative target-only anchor for semantic-skill evaluation \\
RGPE / multisource RGPE & Ranks candidates with source and target surrogate ensembles & Classical transfer-BO comparator for one or multiple source tasks \\
Multitask GP / ICM-BMA & Learns source--target covariance and averages compatible multitask models & Tests whether a numerical transfer model outperforms language-guided transfer \\
Warm-start-only & Keeps source-informed initial target points, then sets sequential transfer mass to zero & Isolates early experimental-design effects \\
Fixed data-only transfer & Uses measured source outcomes with a deterministic role map and no LLM patch & Tests whether the LLM adds value beyond source data and fixed rules \\
\bottomrule
\end{tabularx}
\end{table}

\section{Frozen semantic-skill results}
\label{sec:semantic_results}

One offline model call generates each skill library. Five successful generation
calls in the frozen study used 13,324 prompt tokens and 9,136 completion tokens.
Replay makes zero model calls. Development used 30 calibration seeds and 50
held-out development seeds. Selected skills were then frozen and evaluated on
300 or 500 fresh paired seeds.

\begin{table}[t]
\centering
\caption{Frozen LLM semantic skills versus the strongest target-only anchor.
Intervals are paired normal 95\% confidence intervals over seeds.}
\label{tab:semantic_headline}
\small
\begin{tabularx}{\textwidth}{>{\raggedright\arraybackslash}X >{\raggedright\arraybackslash}X c >{\raggedright\arraybackslash}X c c c}
\toprule
Domain / target & Frozen skill & Seeds & Target-only anchor & $\dfinal$ & $\dauc$ & $\Delta$ top-10 \\
\midrule
Molecular / FreeSolv & Exploration-first kernel skill & 500 & Target portfolio & $+1.433$ & $+1.384$ & $+0.058$ \\
Molecular / Lipophilicity & Shared acquisition schedule & 300 & Target portfolio & $+0.485$ & $+0.607$ & $+0.040$ \\
Reaction / Buchwald--Hartwig & High-molecular-weight ligand prior & 500 & Mixed-kernel GP-EI & $+3.121$ & $+2.859$ & $+0.176$ \\
Materials / experimental band gap & Transition-metal counter prior & 500 & Mixed-kernel GP-EI & $+7.224$ & $+5.930$ & $+0.164$ \\
\bottomrule
\end{tabularx}
\end{table}

The 95\% intervals for AUC were $[0.579,2.189]$ for FreeSolv,
$[0.130,1.084]$ for Lipophilicity, $[1.960,3.757]$ for
Buchwald--Hartwig, and $[4.159,7.701]$ for experimental band gap. These
results establish positive task-specific examples across three scientific
task families. They do not establish that a single skill transfers everywhere.
The ChemLex calibration and held-out check remained negative and was not
promoted to a confirmatory run.

\subsection{Semantic-rule and prior ablations}

Schedule-only keeps the LLM-selected acquisition schedule but removes semantic
rule features. No-prior keeps the rule features but initializes their
coefficient directions at zero. Both controls use the same 500 target seeds as
the full skill.

\begin{table}[t]
\centering
\caption{Component attribution for the two strongest semantic-skill targets.
Rows report full skill minus ablation.}
\label{tab:semantic_ablation}
\begin{tabular}{lrrr}
\toprule
Target / ablation & $\dfinal$ & $\dauc$ & $\Delta$ top-10 \\
\midrule
Buchwald--Hartwig / schedule-only & $+2.745$ & $+2.824$ & $+0.188$ \\
Band gap / schedule-only & $+6.678$ & $+5.065$ & $+0.148$ \\
Buchwald--Hartwig / no-prior & $+1.655$ & $+1.993$ & $+0.094$ \\
Band gap / no-prior & $+2.054$ & $+1.378$ & $+0.034$ \\
\bottomrule
\end{tabular}
\end{table}

The semantic partition and the proposed direction each contribute on these two
targets. A separate causality audit shows a complementary pattern: on Matbench
band gap, the target-only \code{high\_element\_count} feature without an LLM
coefficient prior improved AUC by $+6.167$ ($[4.361,7.972]$); on
Lipophilicity, \code{high\_aromatic\_content} improved AUC by $+1.261$
($[0.891,1.630]$). These controls imply that the most reliable LLM contribution
is often representation choice, while target evidence should estimate the
numeric weight.

\section{Source evidence and transfer mechanism}
\label{sec:mechanism}

\subsection{Source-schema transfer across five routes}

Five frozen routes were evaluated on 300 disjoint target seeds. The
source-schema condition contains source identity, objective, and field
alignment but no source labels or target outcomes.

\begin{table}[t]
\centering
\caption{Source-schema router minus matched target-only semantic router.}
\label{tab:source_schema}
\begin{tabularx}{\textwidth}{X r r r p{2.2cm}}
\toprule
Source $\rightarrow$ target & $\dfinal$ & $\dauc$ & $\Delta$ top-10 & Result \\
\midrule
ESOL $\rightarrow$ Lipophilicity & $+0.602$ & $+1.376$ & $-0.050$ & Positive AUC \\
Lipophilicity $\rightarrow$ FreeSolv & $-17.099$ & $-29.902$ & $-0.187$ & Negative transfer \\
Experimental gap $\rightarrow$ dielectric & $+15.265$ & $+13.085$ & $+0.203$ & Positive transfer \\
Phonon peak $\rightarrow$ dielectric & $+6.789$ & $+4.416$ & $-0.213$ & AUC gain, top-10 cost \\
ChemLex $\rightarrow$ Buchwald--Hartwig & $+3.971$ & $+4.398$ & $+0.193$ & Positive transfer \\
\bottomrule
\end{tabularx}
\end{table}

\begin{figure*}[t]
\centering
\includegraphics[width=0.92\textwidth]{source_evidence_forest.png}
\caption{Paired source-schema effects on five real source--target routes. The
figure is generated from the archived per-seed metrics; it is not an illustrative
or model-generated graphic.}
\label{fig:source_evidence}
\end{figure*}

Shared descriptors are not sufficient. Lipophilicity and hydration free energy
use related molecular representations but have non-interchangeable objective
directions, producing a large negative route. Likewise, the phonon-to-dielectric
route improves mean search reward while reducing extreme top-10 hits. We
therefore treat transferability as metric- and objective-specific.

\subsection{Warm-start and continuous-transfer controls}

Seven source-outcome routes were rerun with two matched controls. Only
ChemLex $\rightarrow$ Buchwald--Hartwig and dielectric $\rightarrow$
experimental band gap passed the new calibration selector. Their large gains
were fully explained by source-informed initial target points.

\begin{table}[t]
\centering
\caption{Mechanism controls. Composite effects are paired held-out differences.
``Post-init.'' compares full transfer with warm-start-only.}
\label{tab:mechanism_controls}
\small
\begin{tabularx}{\textwidth}{>{\raggedright\arraybackslash}X c c c c}
\toprule
Source $\rightarrow$ target & Final vs BO & Warm-start & Post-init. & LLM vs fixed \\
\midrule
ChemLex $\rightarrow$ Buchwald--Hartwig & $+11.452$ & $+22.460$ & $0$ & $0$ \\
Dielectric $\rightarrow$ experimental gap & $+42.461$ & $+85.616$ & $0$ & $0$ \\
Experimental gap $\rightarrow$ dielectric & $+10.017$ & $-0.912$ & $0$ & $+29.424$ \\
Phonons $\rightarrow$ dielectric & $+3.369$ & $+3.361$ & $-1.296$ & $-2.642$ \\
ESOL $\rightarrow$ Lipophilicity & $+1.089$ & $+0.789$ & $-0.145$ & $-0.223$ \\
FreeSolv $\rightarrow$ Lipophilicity & $+2.126$ & $+0.937$ & $0$ & $0$ \\
Lipophilicity $\rightarrow$ FreeSolv & $+31.760$ & $-1.173$ & $0$ & $0$ \\
\bottomrule
\end{tabularx}
\end{table}

The positive $+29.424$ value for experimental gap $\rightarrow$ dielectric
does not indicate improvement beyond warm-start. The full and warm-start-only
trajectories are identical; the number is positive because the fixed data-only
patch is harmful. The defensible mechanism is improved early experimental
coverage on two routes and rejection of unstable transfer on five routes, not
a general post-initialization source-outcome advantage.

\subsection{Source-outcome assignment falsification}
\label{sec:outcome_falsification}

We next asked whether the measured source outcomes carry information through
their association with specific source candidates, rather than merely supplying
an outcome marginal from which an initial design can be selected. For each of
three frozen routes, we held the source features, target initialization, target
seeds, LLM-compiled role map, executable patches, target optimizer, and budget
fixed. We compared the true feature--outcome assignment with 19 assignments
obtained by permuting measured outcomes across the same source candidates. The
audit used 100 held-out target seeds per route and zero sequential transfer mass,
thereby isolating source-informed initial design.

\begin{table}[t]
\centering
\caption{Frozen source-outcome assignment falsification. Differences are paired
best-so-far AUC means with normal 95\% confidence intervals over 100 target
seeds. $p$ is a one-sided empirical randomization test over 19 matched outcome
permutations.}
\label{tab:outcome_falsification}
\small
\begin{tabularx}{\textwidth}{>{\raggedright\arraybackslash}X r r r}
\toprule
Source $\rightarrow$ target & True vs target-only & True vs mean permutation & $p$ \\
\midrule
ChemLex $\rightarrow$ Buchwald--Hartwig
  & $+9.672\;[+7.405,+11.939]$
  & $+20.347\;[+20.347,+20.347]$ & $0.15$ \\
Dielectric $\rightarrow$ experimental gap
  & $+48.954\;[+44.050,+53.857]$
  & $+47.766\;[+47.766,+47.766]$ & $0.10$ \\
Lipophilicity $\rightarrow$ FreeSolv
  & $-1.027\;[-1.228,-0.826]$
  & $+0.451\;[+0.356,+0.545]$ & $0.20$ \\
\bottomrule
\end{tabularx}
\end{table}

\begin{figure*}[t]
\centering
\includegraphics[width=0.98\textwidth]{source_outcome_falsification.png}
\caption{Source-outcome assignment falsification generated directly from the
archived 300 target-seed jobs. Panel A compares the true source assignment with
the matched target-only optimizer. Panel B compares it with the mean of 19
outcome permutations; labels report one-sided empirical randomization $p$
values.}
\label{fig:outcome_falsification}
\end{figure*}

The true source assignment improves over target-only on the reaction and
materials routes, confirming that their source-informed starts improve finite
budget search. However, none of the three correct assignments is extreme under
the matched permutation null ($p=0.10$--$0.20$). The molecular route remains
worse than target-only despite outperforming the mean permutation. The primary
falsification therefore does not establish that the correct source
feature--outcome association contributes information beyond schema and the
initialization mechanism. The positive warm-start results must not be presented
as verified source-outcome learning.

\section{Frozen external task-family stress test}
\label{sec:flip2_external}

We next froze an evaluation on a task family that was absent from controller
development: the FLIP2 Hydrophobic Core protein-stability benchmark. The
official \code{to-P06241} wild-type split contains 24,935 experimentally
measured variants across three protein backbones. Completed P01053 (6,463
variants) and P0A9X9 (8,500 variants) campaigns were the only source tasks;
P06241 (9,972 variants) was the unseen target. The raw split is distributed by
the FLIP2 authors at
\url{https://flip.protein.properties/assets/splits/hydro/to_P06241.csv.gz}.

The public representation uses only the seven randomized F/I/L/M/V core
residues: residue counts, mean side-chain volume, residue diversity, and the
seven-residue core sequence. Target outcomes were unavailable to the LLM and
all policies before each reveal. The split, 100 seeds, three initial target
observations, 12 reveals, source observation counts, kernel settings, and
comparison rules were frozen before execution. The protocol lock records the
canonical config hash, raw-data hash, and hashes of every participating runner.
No P06241 outcome was used for controller or hyperparameter selection.

\begin{table}[t]
\centering
\caption{Frozen FLIP2 external-task evaluation. AUC differences and normal
95\% confidence intervals are paired against target-only GP-UCB over the same
100 target seeds.}
\label{tab:flip2_external}
\small
\begin{tabularx}{\textwidth}{X r r r r}
\toprule
Policy & Mean AUC & $\dauc$ vs target GP & 95\% CI & AUC win rate \\
\midrule
Target-only GP-UCB & 81.949 & --- & --- & --- \\
Multisource RGPE & 64.990 & $-16.959$ & $[-18.407,-15.511]$ & 0.00 \\
Multisource ICM-BMA & 73.816 & $-8.133$ & $[-9.872,-6.394]$ & 0.12 \\
Fixed multisource skill prior & 58.941 & $-23.008$ & $[-25.657,-20.359]$ & 0.00 \\
\bottomrule
\end{tabularx}
\end{table}

\begin{figure*}[t]
\centering
\includegraphics[width=0.98\textwidth]{external_task_family_transfer.png}
\caption{External protein-backbone transfer audit generated directly from the
frozen result files. Panel A reports 100-seed paired effects for three numerical
transfer baselines. Panel B reports one live Opus 4.8 trajectory and therefore
has no sampling confidence interval. The online LLM selected GP rank one and
then abstained, yielding zero online increment over the same-start GP. The full
LLM trajectory was 7.186 AUC below the fixed-design comparator because its
source-informed initial design was worse.}
\label{fig:flip2_external}
\end{figure*}

The initial LLM hypothesis predicted that low aromatic content and high
beta-branched Val/Ile packing would transfer across protein backbones. Its three
target measurements were 28.589, 34.553, and 67.757. The best observation was
instead a low-diversity methionine-rich core. At the first online round, the
proposer and critic explicitly downgraded the original source rule, set
\code{continue\_source\_transfer=false}, and selected target-GP rank one. The
candidate returned 55.005 against an expected 60.0 (absolute error 4.995).
The frozen one-round authority limit then handed all 11 later rounds to target
GP, avoiding 22 nominal proposer--critic calls. Online CARE and the same-start
GP were therefore identical ($\dauc=0$). The full LLM system was worse than the
fixed initial design ($\dauc=-7.186$).

This experiment is positive evidence for auditability and bounded abstention,
not for efficacy on the new domain. It also demonstrates that a semantically
plausible cross-backbone hypothesis can fail even when source and target share
an apparently identical seven-residue design schema. Any generalization claim
must therefore include negative-transfer rate and same-start target-only
controls.

\subsection{Source-only family gate}
\label{sec:flip2_source_only_gate}

We next asked whether the completed protein campaigns could reject harmful
transfer without using any P06241 outcome. We treated P01053 and P0A9X9 in turn
as pseudo-targets and evaluated the three fixed transfer policies over 50 paired
seeds per direction. A policy was eligible for deployment only if the lower
bound of its paired 95\% confidence interval exceeded zero on \emph{both}
source-only routes. Otherwise the frozen fallback was target-only GP-UCB. The
P06241 target was prohibited from the calibration-route specification and from
all policy-selection calculations.

RGPE changed AUC by $-6.743$ [$-8.466,-5.020$] on P01053$\rightarrow$P0A9X9
and $-5.573$ [$-11.073,-0.073$] on the reverse route. ICM-BMA changed AUC by
$-3.361$ [$-4.739,-1.984$] and $-8.158$ [$-13.936,-2.380$], respectively.
The fixed skill prior changed AUC by $-12.702$ [$-14.327,-11.078$] and
$-27.288$ [$-32.522,-22.054$]. No transfer policy was therefore eligible.
Without the gate, source-only selection would have chosen ICM-BMA because it
had the least negative route-equal mean ($-5.760$). On P06241, that policy lost
$8.133$ AUC relative to target-only GP-UCB; the gate deployed target-only GP-UCB
and therefore avoided that loss by construction.

\begin{figure*}[t]
\centering
\includegraphics[width=0.98\textwidth]{source_only_family_gate.png}
\caption{Source-only family-gate audit generated directly from the frozen
result files. The first two columns are the only calibration evidence: each
completed source backbone is treated as a pseudo-target in turn. All three
transfer policies have confidence intervals below zero on both directions and
are rejected. The third column is an external deployment audit on P06241 and
does not enter calibration or policy selection. The gate falls back to
target-only GP-UCB and avoids the 8.133 AUC loss of the otherwise selected
ICM-BMA policy. The gate design followed observation of the external-family
result, so this is retrospective mechanism evidence rather than prospective
safety confirmation.}
\label{fig:flip2_source_only_gate}
\end{figure*}

This result narrows the next confirmatory question. The numerical rule can be
computed without deployment-target labels and would have blocked the observed
negative transfer, but its design history prevents a prospective claim. This
requirement therefore defined the confirmatory test: preregister the same eligibility
rule, source-only calibration procedure, and target-only fallback on a second
external family before any target outcome is inspected.

\subsection{Preregistered FLIP2 Rhodopsin confirmation}
\label{sec:flip2_rhomax_gate}

We tested that requirement on the official FLIP2 Rhodopsin peak-wavelength
benchmark (Rhomax). Its by-wild-type split contains 884 measured sequences from
75 wild types. The official partition assigns the five most common wild types to
training, 34 wild types to validation, and 36 wild types to test; these contain
584, 116, and 184 sequence variants, respectively. The file is distributed at
\url{https://flip.protein.properties/assets/splits/rhomax/by_wild_type.csv.gz}.
Public candidate features contain only sequence length, amino-acid composition,
and six predeclared physicochemical fractions. The measured peak wavelength is
hidden until reveal and is mapped to the declared replay score by
$100(\lambda-350)/300$, clipped to $[0,100]$.

The protocol, eligibility rule, source-only route construction, target-only
fallback, implementation hashes, and all 200 evaluation seeds were committed as
\code{4130dd4} before the raw file was downloaded. The preregistration record
states that the raw dataset was absent and deployment-target outcomes were
unavailable at freeze. Calibration used train$\rightarrow$validation and
validation$\rightarrow$train, with 50 paired seeds per route. The official test
partition was excluded from gate computation and method selection. A method was
eligible only if its paired AUC 95\% confidence-interval lower bound exceeded
zero on both calibration routes.

\begin{table}[t]
\centering
\caption{Preregistered FLIP2 Rhodopsin family-gate confirmation. Values are
paired $\dauc$ relative to target-only GP-UCB. The official-test column contains
100 paired seeds and was read only after the frozen gate decision.}
\label{tab:flip2_rhomax_gate}
\scriptsize
\setlength{\tabcolsep}{2.5pt}
\renewcommand{\arraystretch}{0.92}
\begin{tabularx}{\textwidth}{@{}>{\raggedright\arraybackslash}p{2.7cm} >{\centering\arraybackslash}X >{\centering\arraybackslash}X >{\centering\arraybackslash}X@{}}
\toprule
Route & RGPE & ICM-BMA & Fixed skill prior \\
\midrule
Train$\rightarrow$validation & $-0.298$ [$-0.922,+0.325$] & $-0.791$ [$-1.397,-0.185$] & $+0.384$ [$-0.301,+1.069$] \\
Validation$\rightarrow$train & $-0.263$ [$-1.252,+0.727$] & $-0.965$ [$-1.660,-0.270$] & $+1.192$ [$+0.271,+2.112$] \\
Official test & $-3.819$ [$-4.498,-3.140$] & $-1.539$ [$-2.066,-1.011$] & $-1.548$ [$-1.996,-1.099$] \\
\bottomrule
\end{tabularx}
\end{table}

The route-equal calibration mean was largest for the fixed skill prior
($+0.788$), but its confidence interval crossed zero on
train$\rightarrow$validation, so it was ineligible. The gate therefore deployed
target-only GP-UCB. After deployment was fixed, the official test showed that
the ungated choice would have lost 1.548 AUC; RGPE and ICM-BMA were also harmful.
The target-only fallback therefore avoided 1.548 AUC of negative transfer.

\begin{figure*}[t]
\centering
\includegraphics[width=0.98\textwidth]{prospective_rhomax_gate.png}
\caption{Prospective external-family gate confirmation generated from the
preregistered result files. Panel A shows paired AUC effects and 95\% confidence
intervals on the two source-only calibration routes and the untouched official
test partition. Panel B records the predeclared decision: no transfer method met
the requirement that the confidence-interval lower bound exceed zero on both
calibration routes, so the policy fell back to target-only GP-UCB. Official test
outcomes enter only the post-decision evaluation.}
\label{fig:flip2_rhomax_gate}
\end{figure*}

This result closes the prospective-design gap for one safety endpoint: a rule
frozen before data access prevented a statistically stable negative transfer on
a second external protein family. It does not show that CARE improves Rhodopsin
optimization, that every harmful transfer will be rejected, or that the system
generalizes positively beyond replay. The next subsection tests positive
external efficacy separately; a paired prospective laboratory campaign remains
an open requirement.

\subsection{Preregistered FLIP2 IRED positive-transfer confirmation}
\label{sec:flip2_ired_positive}

The preceding experiments test whether CARE can refuse an unsafe transfer. We
next preregistered a distinct positive-efficacy question on the official FLIP2
Imine Reductase (IRED) \code{two-to-many} split. IREDs catalyze the reduction of
imines to amines. The official training and validation partitions contain 3,746
and 662 measured variants with zero, one, or two mutations; the official test
partition contains 4,178 variants with three to fifteen mutations. The raw file
is distributed at
\url{https://flip.protein.properties/assets/splits/ired/two_to_many.csv.gz}.

Before that file was downloaded, Claude Opus 4.8 received only the public task
description, partition counts, scientific context, optimization endpoint, and
a two-item executable menu: a bounded additive-mutation skill or abstention to
target-only GP-UCB. It proposed the falsifiable hypothesis that shrunk effects
estimated from low-order variants would remain sufficiently additive to rank
higher-order variants. It also recorded four failure conditions: dominant
epistasis, unseen mutation positions, low-to-high-order distribution shift, and
excessive early source authority. The model prompt, raw response, parsed
hypothesis, token count, and request hash were archived before data access.

The resulting skill estimates exact mutation-token effects from source
observations, backs off to amino-acid physicochemical class transitions when an
exact token is unavailable, sums the available contributions for each target
variant, rank-normalizes the result, and blends it with target GP-UCB. Source
authority decays from 0.50 to 0.10 over the twelve target reveals. The skill
never reads an official-test activity during construction. The executable
skill, candidate menu, calibration rule, code hashes, 100 calibration seeds,
100 official-test seeds, three initial target observations, twelve reveals,
and no-optional-stopping rule were committed as \code{5389a9c} while the raw
dataset was absent.

Calibration used only official train$\rightarrow$validation replay. The
predeclared LLM skill was eligible only if its paired AUC normal 95\% confidence
interval had a lower bound above zero. It passed with $+0.468$
[$+0.365,+0.571$]. The gate then deployed that unchanged skill from
train$+$validation to the untouched official test. Table
\ref{tab:flip2_ired_positive} reports paired-bootstrap intervals with 200,000
resamples as a sensitivity analysis; they agree with the preregistered normal
interval decision.

\begin{table}[t]
\centering
\caption{Preregistered FLIP2 IRED positive-transfer confirmation. Values are
paired $\dauc$ against target-only GP-UCB. W/T/L is computed on the 100 official
test seeds. Only the LLM additive skill was a gate candidate; the other methods
are predeclared comparators and are shown to avoid a weak-baseline claim.}
\label{tab:flip2_ired_positive}
\scriptsize
\setlength{\tabcolsep}{2.5pt}
\renewcommand{\arraystretch}{0.94}
\begin{tabularx}{\textwidth}{@{}X r r r c@{}}
\toprule
Method & Train$\rightarrow$val. [95\% CI] & Official test [95\% CI] & Test W/T/L & Gate candidate \\
\midrule
RGPE & $+0.137$ [$+0.068,+0.209$] & $-0.085$ [$-0.207,+0.032$] & 44/13/43 & No \\
ICM-BMA & $+0.032$ [$-0.032,+0.098$] & $-0.194$ [$-0.302,-0.091$] & 23/17/60 & No \\
Fixed skill prior & $+0.489$ [$+0.385,+0.595$] & $+0.864$ [$+0.709,+1.019$] & 78/8/14 & No \\
LLM additive skill & $+0.468$ [$+0.366,+0.574$] & $+0.389$ [$+0.274,+0.503$] & 72/9/19 & Yes \\
\bottomrule
\end{tabularx}
\end{table}

\begin{figure*}[t]
\centering
\includegraphics[width=0.98\textwidth]{prospective_ired_positive_transfer.png}
\caption{Prospective IRED positive-transfer result generated directly from the
paired metrics and 6,000 official-test audit rows. Panel A shows paired AUC
effects with 200,000-resample bootstrap intervals on calibration and untouched
test seeds. Panel B shows mean best-so-far activity after each reveal; ribbons
are bootstrap 95\% intervals over the same 100 seeds. The fixed prior is shown
because it is a strong predeclared comparator, not an LLM ablation.}
\label{fig:flip2_ired_positive}
\end{figure*}

On the official test, the deployed LLM skill improved AUC by 0.389
[$0.274,0.503$], with 72 wins, nine ties, and 19 losses; the two-sided exact
sign-test value excluding ties was $p=1.97\times10^{-8}$. Its final-best effect
was $+0.162$ [$-0.066,+0.390$], so the evidence is specifically earlier discovery
under the predeclared AUC endpoint, not a stable terminal-best improvement. The
top-10 hit-rate difference was $-0.13$ [$-0.207,-0.053$]. The fixed skill prior
was substantially stronger in AUC ($+0.864$), demonstrating that the
outcome-blind LLM hypothesis was useful but not uniquely or maximally useful.

All 1,200 LLM-skill decision rows are retained. In illustrative seed 97200, the
skill began with source mass 0.50, 96.1\% exact mutation-token coverage, 3.9\%
class-transition fallback, and complete candidate evidence coverage. By round
12, source mass had decayed to 0.10 and the best observed activity was 1.355.
That seed's AUC was 0.463 versus $-0.243$ for matched target-only GP-UCB. This
trace shows the declared mechanism, but the claim rests on all 100 paired seeds,
not on the selected example. The experiment supplies prospective positive
finite-pool transfer evidence in one external enzyme family. It does not prove
universal cross-domain generalization, LLM superiority over fixed priors, or
prospective wet-lab discovery.

\subsection{Post-hoc IRED source-outcome assignment falsification}
\label{sec:ired_outcome_falsification}

The prospective IRED endpoint establishes that the deployed additive skill
improves early finite-pool search. It does not by itself establish why. We
therefore performed a stricter post-hoc mechanism audit that asks whether the
measured mutation--activity binding is more useful than false bindings. The
official test candidates, 100 target seeds, three initial observations, twelve
reveals, GP kernel, additive executor, 0.50-to-0.10 source-authority schedule,
3,500 source feature rows, and source outcome marginal distribution were held
fixed. The only intervention was the assignment of measured source outcomes to
source variants. We compared the true assignment with 99 deterministic outcome
permutations.

\begin{table}[t]
\centering
\caption{IRED source-outcome assignment falsification. Intervals resample the
100 paired target seeds. The assignment-level $p$-value instead ranks the true
binding among the 99 outcome-permuted bindings and uses a one-sided plus-one
correction. The audit was designed after the prospective IRED result and is
therefore a mechanism analysis, not independent confirmation.}
\label{tab:ired_outcome_falsification}
\scriptsize
\setlength{\tabcolsep}{3.2pt}
\begin{tabularx}{\textwidth}{@{}X r r r@{}}
\toprule
Comparison & Mean $\dauc$ & 95\% target-seed interval & Assignment result \\
\midrule
True binding vs. target GP-UCB & $+0.389$ & [$+0.274,+0.504$] & 72/9/19 W/T/L \\
Permutation mean vs. target GP-UCB & $+0.049$ & [$-0.037,+0.132$] & 46/99 means $>0$ \\
True binding vs. permutation mean & $+0.340$ & [$+0.259,+0.422$] & $p=0.07$; 6 nulls $\geq$ true \\
\bottomrule
\end{tabularx}
\end{table}

\begin{figure*}[t]
\centering
\includegraphics[width=0.98\textwidth]{ired_source_outcome_falsification.png}
\caption{IRED source-outcome mechanism audit generated directly from 10,000
skill trajectories, 100 matched target-GP trajectories, and the archived
assignment schedule. Panel A locates the true skill effect within the 99
outcome-permuted assignment effects. Panel B compares the mean permuted skill,
the LLM-selected additive skill with the true binding, and the predeclared fixed
rank prior against target-only GP-UCB. Panel C shows best-so-far trajectories
under the same target budget. Ribbons and effect bars are bootstrap 95\%
intervals over target seeds.}
\label{fig:ired_outcome_falsification}
\end{figure*}

The true assignment exceeds the mean outcome-permuted skill by 0.340 AUC with a
target-seed bootstrap interval entirely above zero. This answers a conditional
question: for the declared set of 99 false bindings, the measured binding is
better on average. The finite assignment randomization test asks a different
question: is the true assignment sufficiently extreme among independently
permuted assignments? Six null assignments achieved an equal or larger mean
AUC effect, giving $p=(1+6)/(1+99)=0.07$. Thirty null assignments also had
target-seed intervals above zero against target GP-UCB. Consequently, this audit
provides suggestive but non-unique mechanism evidence. It does not prove that
the LLM identified the uniquely correct source feature--outcome relation. The
next external protocol should freeze this assignment null before outcome access
and treat mechanism attribution as a co-primary endpoint.

\subsection{Preregistered FLIP2 TrpB efficacy and source-outcome confirmation}
\label{sec:flip2_trpb_replication}

We implemented that next protocol on a previously unused external family: the
official FLIP2 Tryptophan Synthase $\beta$-Subunit (TrpB)
\code{one-to-many} split. The public benchmark contains 228,298 measured
variants across ten combinatorial sub-landscapes and 20 engineered positions.
The official train and validation partitions contain 304 and 76 zero- or
single-mutation variants; the official test contains 227,918 variants with two
to four mutations. Before downloading the file, we froze an outcome-blind
finite pool consisting of at most 5,000 rows per partition selected solely by
the lexicographically smallest SHA-256 sequence hashes. The resulting test
pool contains 215 double, 1,899 triple, and 2,886 quadruple mutants.

Claude Opus 4.8 received only the public task description, official split
logic, scientific context, budget, finite-pool rule, and a two-item menu: the
same bounded additive-mutation executor or abstention. It selected the
additive skill and predicted that first-order mutation effects would improve
early target ranking while explicitly naming dominant epistasis, incomplete
mutation coverage, low-to-high-order shift, and authority miscalibration as
failure conditions. The raw response, 2,460-token usage record, parsed
hypothesis, request hash, executable code hashes, 100 calibration seeds, 100
deployment seeds, and a 99-assignment source-outcome randomization test were
committed as \code{c5f4065} while the raw TrpB file was absent.

Table \ref{tab:flip2_trpb_replication} separates candidate-skill efficacy from
the frozen routing decision. On the official test, the LLM additive skill
improved best-so-far AUC by 0.949 [95\% paired CI $0.864,1.034$] with 97 wins,
three ties, and no losses. The generic skill prior improved by 0.154, whereas
RGPE and ICM-BMA did not improve. This comparison rules out a weak target-only
baseline explanation and shows that mutation-specific source outcomes carry
substantially more useful information than a generic source ranking.

\begin{table}[t]
\centering
\caption{Preregistered FLIP2 TrpB external replication. Official-test values
are paired best-so-far AUC differences against target-only GP-UCB over 100
seeds. The gate decision was fixed before those outcomes were inspected.}
\label{tab:flip2_trpb_replication}
\scriptsize
\setlength{\tabcolsep}{3.0pt}
\begin{tabularx}{\textwidth}{@{}X r r r@{}}
\toprule
Method or decision & Mean $\dauc$ & 95\% paired CI & W/T/L or result \\
\midrule
RGPE & $-0.023$ & [$-0.057,+0.011$] & 29/19/52 \\
ICM-BMA & $-0.037$ & [$-0.062,-0.011$] & 22/25/53 \\
Generic skill prior & $+0.154$ & [$+0.091,+0.216$] & 64/13/23 \\
LLM additive skill & $+0.949$ & [$+0.864,+1.034$] & 97/3/0 \\
Frozen deployed policy & $0.000$ & exact target-only fallback & HOLD \\
\bottomrule
\end{tabularx}
\end{table}

The predeclared assignment test then held the target pool, 100 new target
seeds, starting observations, budget, GP kernel, additive executor, source
rows, outcome marginal distribution, and 0.50-to-0.10 authority schedule fixed.
Only the binding between 380 source outcomes and source mutation features was
changed. With the measured binding, the skill improved AUC by 0.885 [normal
95\% CI $0.805,0.965$] versus target-only GP-UCB. Relative to the per-seed mean
of 99 outcome-permuted skills, the true binding improved AUC by 0.890
[200,000-resample bootstrap CI $0.810,0.970$]. No null assignment reached the
true assignment effect, giving the predeclared plus-one randomization value
$p=(1+0)/(1+99)=0.01$. The mechanism co-primary endpoint therefore passed.

\begin{figure*}[t]
\centering
\includegraphics[width=0.99\textwidth]{prospective_trpb_replication.png}
\caption{Preregistered TrpB replication generated from the complete real-data
result files. Panel A compares the LLM additive skill with strong matched
baselines on the untouched sequence-hash test pool. Panel B locates the true
source binding among 99 outcome-permuted assignments. Panel C shows matched
best-so-far trajectories under the same target budget. The red annotation
records the frozen gate's HOLD decision rather than silently replacing it with
the successful ungated candidate.}
\label{fig:flip2_trpb_replication}
\end{figure*}

The gate result reveals a routing failure rather than an efficacy failure. Its
train$\rightarrow$validation estimate was $+0.044$ [95\% CI
$-0.018,+0.106$], so the frozen lower-bound rule deployed target-only GP-UCB.
That calibration route had zero exact mutation-token coverage because both
partitions contain held-out single mutants and the executor could use only
class-level fallback. At official-test deployment, the source train+validation
set covered 99.97\% of mutation tokens in the multi-mutant target pool. The
calibration task therefore tested a different information regime from the
deployment task. We retain the abstention as the complete-system result: the
LLM proposed a strongly useful skill and the attribution test confirmed its
source dependence, but the current router failed to deploy it. A revised gate
must condition on structural coverage and must be frozen on another family;
the TrpB outcome cannot be used to retroactively change this decision.

\subsection{Cross-family prospective synthesis}
\label{sec:external_triad}

The three preregistered external-family experiments answer different questions
and should not be collapsed into a single efficacy average. Rhomax tests
prospective safety: the frozen gate rejected a harmful fixed prior and kept the
complete router at the target-only reference. IRED tests positive deployment:
the LLM-selected additive skill passed calibration and the complete router
realized a $+0.389$ AUC effect. TrpB jointly tests efficacy and source-outcome
attribution: the same skill class improved AUC by $+0.949$ and passed the
predeclared assignment test, but the complete router remained at zero because
the gate issued a false-negative hold.

\begin{table}[t]
\centering
\caption{Descriptive synthesis of preregistered external-family results. Each
deployment comparison contains 100 paired target seeds. Candidate efficacy and
the complete deployed router are reported separately. IRED uses its
preregistered paired-bootstrap interval; Rhomax and TrpB use paired normal
intervals from their frozen protocols.}
\label{tab:external_triad}
\scriptsize
\setlength{\tabcolsep}{3.0pt}
\begin{tabularx}{\textwidth}{@{}l X r l r X@{}}
\toprule
Family & Candidate & Candidate $\dauc$ [95\% CI] & Gate & Deployed $\dauc$ & Router outcome \\
\midrule
Rhomax & Fixed skill prior & $-1.548$ [$-1.996,-1.099$] & Hold & $0.000$ & Correct rejection \\
IRED & LLM additive skill & $+0.389$ [$+0.274,+0.502$] & Transfer & $+0.389$ & Correct deployment \\
TrpB & LLM additive skill & $+0.949$ [$+0.864,+1.034$] & Hold & $0.000$ & False-negative hold \\
\bottomrule
\end{tabularx}
\end{table}

\begin{figure*}[t]
\centering
\includegraphics[width=0.99\textwidth]{prospective_external_family_triad.png}
\caption{Candidate efficacy and complete-router behavior across three
preregistered external protein families. Circles and paired 95\% confidence
intervals show the candidate policy relative to target-only GP-UCB; diamonds
show the policy actually deployed by the frozen router. The plot is generated
directly from the archived result tables. It is a descriptive synthesis, not a
population estimate of generalization because only three families are present.}
\label{fig:external_triad}
\end{figure*}

This synthesis identifies a routing objective that a success-only aggregate
would hide. The candidate layer distinguishes useful from harmful transferable
structure in these examples, but the router has one false negative among two
beneficial candidates. The next protocol therefore freezes a coverage-aware
decision rule on a previously unused Alpha Amylase family. That experiment is
prospective only if its public task specification, LLM hypothesis, thresholds,
code hashes, and target-outcome exclusion are committed before the official file
is downloaded. The existing TrpB outcomes are development evidence for this new
router and cannot contribute to its confirmatory endpoint.

\section{Online LLM controller and strong baselines}
\label{sec:online}

\subsection{Frozen route portfolio}

The primary online comparison is the LLM controller versus the same LLM initial
target observations followed by target-only GP-UCB. In the frozen 11-route
first-trajectory portfolio, six routes improve AUC, two tie, and three lose.
The equal-route mean is $+1.753$ with a route-bootstrap 95\% interval
$[0.415,3.315]$. Because each route contains one stochastic LLM trajectory and
some routes share task families, this interval describes route variation, not
within-route model-call uncertainty.

\begin{table}[t]
\centering
\caption{Aggregate online-controller and handoff audits relative to same-start
target GP-UCB. Controller rows are retrospective replays.}
\label{tab:controller_results}
\begin{tabularx}{\textwidth}{X r r p{2.6cm} r}
\toprule
Evidence tier / method & Routes & Mean $\dauc$ & 95\% route-bootstrap interval & Win/tie/loss \\
\midrule
Frozen online LLM portfolio & 11 & $+1.753$ & $[+0.415,+3.315]$ & $6/2/3$ \\
Cross-validated round-three gate & 11 & $+2.036$ & $[+0.425,+3.901]$ & $5/4/2$ \\
Fixed threshold-5 gate replay & 11 & $+2.135$ & $[+0.515,+4.001]$ & $6/3/2$ \\
Complete retrospective inventory & 17 & $+0.762$ & $[-0.501,+2.091]$ & $8/4/5$ \\
Route-disjoint full online LLM & 6 & $-1.056$ & $[-3.010,+0.154]$ & $2/2/2$ \\
Route-disjoint threshold-5 gate & 6 & $-1.968$ & $[-6.438,+0.873]$ & $2/2/2$ \\
Route-disjoint one-round authority & 6 & $+0.612$ & $[-0.442,+1.793]$ & $2/3/1$ \\
\bottomrule
\end{tabularx}
\end{table}

The 17-route inventory includes earlier null and negative development routes
and is reported to prevent route selection by outcome. The round-three gate
improves the original development panel but fails on six route-disjoint routes.
A training-only comparison of 115 handoff policies selects one-round bounded
authority; its disjoint interval still crosses zero. In a leave-one-route-out
stability audit over the 11 development routes, the same one-round policy was
re-selected in 10 of 11 folds. The policy selected in each fold achieved a
held-out mean $\dauc=+1.141$ with $2/8/1$ wins/ties/losses. This supports
selection stability, but the remaining loss and retrospective design preclude
a universal safe-transfer claim.

\begin{figure*}[t]
\centering
\includegraphics[width=0.96\textwidth]{predeclared_baseline_portfolio.png}
\caption{Frozen route portfolio relative to predeclared target GP-UCB and
classical transfer-BO comparators. Baseline rules were fixed at analysis time;
the strongest-realized comparator is shown only as a post-hoc stress test.}
\label{fig:baseline_portfolio}
\end{figure*}

\begin{figure*}[t]
\centering
\includegraphics[width=0.98\textwidth]{controller_selection_stability.png}
\caption{Controller selection stability and route-disjoint effects. The
one-round bounded-authority policy is re-selected in 10 of 11
leave-one-development-route-out folds (left). Frozen application to six
disjoint retrospective trajectories yields two wins, three ties, and one loss
against same-start target-only GP-UCB (right).}
\label{fig:bounded_authority}
\end{figure*}

\subsection{Archived probability calibration and abstention implication}

The final proposer--critic response reports a probability that the executed
candidate will improve the pre-round best. We scored this field on every one of
the 170 archived decisions and compared it with the production GP probability
of improvement for the same executed candidate. This design evaluates
probability estimation without confounding it with candidate selection. There
were 27 improvements, an observed rate of 0.159.

\begin{table}[t]
\centering
\caption{Retrospective probability calibration on the complete 17-route
archive. ECE uses five fixed 0.2-wide bins.}
\label{tab:probability_calibration}
\begin{tabular}{lrrrrrr}
\toprule
Probability source & Mean $p$ & Observed & Bias & Brier & ECE & ROC AUC \\
\midrule
Final LLM decision & 0.219 & 0.159 & $+0.060$ & 0.133 & 0.060 & 0.679 \\
GP, same candidate & 0.182 & 0.159 & $+0.023$ & 0.134 & 0.023 & 0.628 \\
\bottomrule
\end{tabular}
\end{table}

The equal-route mean LLM-minus-GP Brier difference is $-0.0016$ with a
route-bootstrap 95\% interval $[-0.0132,+0.0090]$; negative values favor the
LLM. The corresponding equal-route ROC-AUC difference is $+0.0308$
$[-0.0581,+0.1373]$. Both intervals include zero. The LLM probability is on
average six percentage points above the realized improvement rate, whereas the
GP probability is 2.3 points high. Thus, the archive does not establish that
the LLM is better calibrated, and its self-reported probability should not be
used directly as a deployment gate.

\begin{figure*}[t]
\centering
\includegraphics[width=0.98\textwidth]{llm_probability_calibration.png}
\caption{Probability calibration on real archived decisions. Panel A compares
fixed-bin predicted and observed improvement rates. Panel B shows the
route-level Brier difference; negative values favor the LLM. Of the 170 GP
probabilities, 110 were archived directly and 60 were reconstructed from the
archived posterior mean, standard deviation, and current best using the
production Gaussian probability-of-improvement formula. The audit pools two
model/controller versions and is not a prospective gate validation.}
\label{fig:probability_calibration}
\end{figure*}

The operational implication is to fit a probability calibration map only on
declared development routes, use the calibrated probability as an abstention
signal, and then freeze both map and threshold before disjoint evaluation. This
prospective test remains outstanding.

\subsection{Live bounded-authority pilot}

One real Opus 4.8 proposer--critic decision completed on each of six disjoint
routes. All six selected the same candidate as GP-UCB rank one. The online AUC
increment was therefore exactly zero for every route, with no GP override and
no critic revision. Three traces retained the source-transfer hypothesis in
text, but none changed the executed action. The six calls consumed 99,243
tokens. Returning control after one reveal avoided 108 nominal later
proposer--critic calls. One pre-response credential failure is preserved and is
terminal under the frozen retry rule; therefore the launch cannot satisfy the
declared 180-trajectory confirmation criterion.

\begin{figure*}[t]
\centering
\includegraphics[width=0.95\textwidth]{decision_impact.png}
\caption{Decision-impact audit for the live bounded-authority pilot. Structured
LLM reasoning was produced, but all executed candidates matched GP rank one.
The distinction between participation and causal action change is central to
the evaluation.}
\label{fig:decision_impact}
\end{figure*}

\section{Representative audited trajectories}
\label{sec:cases}

The following cases are post-hoc mechanism audits, not independent statistical
evidence. At every round the target value of an unexecuted candidate remained
hidden until selection.

\begin{table}[t]
\centering
\caption{Representative online decisions relative to the same-start target GP.}
\label{tab:cases}
\begin{tabularx}{\textwidth}{p{3.4cm}p{2.4cm}X r r}
\toprule
Source $\rightarrow$ target & Domain & Audited decision & $\dauc$ & $\dfinal$ \\
\midrule
FreeSolv $\rightarrow$ Lipophilicity & Molecular & GP rank 4 selected in round 1 & $+6.325$ & $+4.875$ \\
Experimental gap $\rightarrow$ dielectric & Materials & GP rank 5 tellurium candidate selected in round 1 & $+6.745$ & $0.000$ \\
Suzuki cases 1--3 $\rightarrow$ case 4 & Reaction & High-temperature catalyst hypothesis retained & $+1.380$ & $+5.700$ \\
Phonon peak $\rightarrow$ bulk modulus & Materials & Direct source relation rejected; target search rerouted & $+3.240$ & $+4.324$ \\
\bottomrule
\end{tabularx}
\end{table}

\noindent\textbf{FreeSolv to Lipophilicity.}
The initial target best was 67.25. Target GP ranked \code{lipo\_2256} first;
the LLM instead selected \code{lipo\_0635}, GP rank 4 with nearly identical
expected improvement. It cited a more halogenated, less polar aromatic
hypothesis. The reveal was 86.875, raising the best value by 19.625 in round
one. CARE finished with AUC 89.575 and final best 89.875, versus 83.250 and
85.000 for the matched GP.

\noindent\textbf{Experimental band gap to dielectric response.}
The transferred hypothesis used the inverse relation between band gap and
polarizability to favour heavy chalcogenides. In round one, the controller
selected \code{real\_matbench\_dielectric\_04424}, GP rank 5. Its hidden value
was 59.5114, the trajectory maximum. Both policies eventually reached the same
final value, but CARE reached it earlier, producing $\dauc=+6.745$.

\noindent\textbf{Suzuki cases 1--3 to case 4.}
Source campaigns supported a high-temperature catalyst hypothesis; temperature
had the largest recorded source correlation (0.647). The initial target set
contained a 94.1-yield P1--L1 condition at 110\,$^\circ$C. The controller
retained high-temperature catalyst variants and later revealed a 99.8-yield
condition, while the same-start GP finished at 94.1.

\noindent\textbf{Phonon peak to bulk modulus.}
The controller set \code{continue\_source\_transfer=false} in every online
round after finding no reliable phonon-to-modulus ranking. It nevertheless
used composition semantics and target history to change four GP choices,
finishing at 77.6127 versus 73.2886. This is evidence for abstention and target
adaptation, not direct transfer of phonon outcomes.

\section{Why some LLM skills help and others fail}
\label{sec:why}

The ablations support three necessary conditions rather than a generic LLM
advantage. First, the semantic partition must align with the target property.
Element-count and aromaticity partitions are positive examples. Second, the
skill must change early sampling or a later executed action. The largest
source-conditioned mechanism-control effects come from initial design; the
six-route live pilot has zero online increment because no action changes.
Third, target evidence must recalibrate the magnitude. No-prior controls often
retain or improve performance, showing that the LLM is more reliable at
proposing a representation than a numeric effect size.

Failure modes follow the same logic. Lipophilicity $\rightarrow$ FreeSolv
shares molecular descriptors but not an interchangeable objective direction,
producing $\dauc=-29.902$. Source outcome summaries never improve over
source-schema-only skills in the four causality pairs. The current prompt is
therefore not a validated mechanism for transporting empirical effect sizes.
The operational decision rule is: deploy only when semantic alignment,
action change, and target recalibration are all observable; otherwise hold,
freeze, or return to the matched target-only optimizer.

\section{Statistical analysis and confirmation protocol}
\label{sec:statistics}

Semantic-skill tables report paired seed differences and normal 95\%
confidence intervals. Online portfolio intervals resample routes with equal
route weight. Two-sided sign tests are secondary. These route intervals do not
replace repeated stochastic LLM calls, and the older semantic-study intervals
are not adjusted for multiple comparisons.

The prospective IRED gate used its preregistered normal paired interval for the
binary deployment decision. For reporting, Table
\ref{tab:flip2_ired_positive} additionally gives deterministic percentile
paired-bootstrap intervals with 200,000 resamples and exact two-sided sign
tests after excluding ties. All 100 declared calibration seeds and all 100
declared official-test seeds are included. The bootstrap analysis was added
after execution as a robustness summary and did not alter the frozen gate.

The frozen repeated protocol declares 30 independent trajectories per route,
the model, temperature, prompt, candidate menu, target budget, retry semantics,
same-start GP comparator, and primary hierarchical-bootstrap analysis. No
confirmatory interval is reported until all declared trajectories complete.
API cost-cap, rate-limit, timeout, and transport failures are retained as
infrastructure events and do not count as negative scientific outcomes. A poor
reward never triggers a retry. Route-level secondary tests use
Benjamini--Hochberg correction.

The first repeated trajectory on each of 11 routes reproduced the aggregate
count of six wins, two ties, and three losses, but four individual routes
reversed sign relative to the earlier first trajectories. The mean absolute
paired AUC change was 1.048. This descriptive stability audit motivates the
30-trajectory protocol and prevents a single stochastic run from defining a
route as transferable.

\section{Datasets, metrics, and reproducibility}
\label{sec:repro}

Experiments use public measured finite pools from Dreher--Doyle
Buchwald--Hartwig HTE, ChemLex Acid--Amine wet-lab data, Reizman Suzuki cases,
MoleculeNet ESOL, FreeSolv, Lipophilicity and BACE, Matbench experimental band
gap, dielectric, phonons, bulk modulus, shear modulus, perovskites and steels,
Materials Project-derived property pools, and the official FLIP2 Hydrophobic
Core, Rhodopsin, and Imine Reductase splits. Synthetic Suzuki and synthetic
ChemLex adapters are retained only as smoke tests and are not used for the
headline real-data claims.

Every reported run stores the config, config hash, seed ranges, target budget,
candidate IDs, selected outcomes, per-round best-so-far trace, model request
and response, validation decision, fallback, token usage, and SHA-256 sidecars.
The code and frozen artifacts are available at
\url{https://github.com/SHITIANYU-hue/CARE2.0/tree/exploration-aware-llm}.
Credentials are never stored in the repository.

\subsection{Primary artifact map}

\begin{table}[!htbp]
\centering
\caption{Repository artifacts supporting the supplement. Paths are relative to
\code{experiments/care\_replay/results}.}
\label{tab:artifact_map}
\small
\begin{tabularx}{\textwidth}{>{\raggedright\arraybackslash}p{6.2cm}X}
\toprule
Artifact & Contents \\
\midrule
2026-07-21: cross-domain semantic skills & Frozen 300/500-seed semantic-skill results, ablations, calls, and traces \\
2026-07-21: LLM evidence causality & Full, source-schema-only, target-only, schedule-only, and no-prior conditions \\
2026-07-23: source-evidence extension & Five 300-seed source-schema transfer routes and round-efficiency analysis \\
2026-08-07: mechanism controls & Warm-start-only and fixed-data-only controls on seven routes \\
2026-08-15: Opus generalization study & Frozen online LLM route portfolio and complete call traces \\
2026-08-24: predeclared baseline portfolio v2 & Target GP, RGPE, and multitask-GP matched baselines \\
2026-08-24: route-disjoint bounded authority & Training/evaluation route split and 115-policy handoff audit \\
2026-08-24: controller selection stability & Eleven leave-one-development-route-out re-selections and six frozen disjoint-route effects \\
\bottomrule
\end{tabularx}
\end{table}

\begin{table}[!htbp]
\centering
\caption{Repository artifacts supporting the supplement (continued).}
\label{tab:artifact_map_continued}
\small
\begin{tabularx}{\textwidth}{>{\raggedright\arraybackslash}p{6.2cm}X}
\toprule
Artifact & Contents \\
\midrule
2026-08-24: bounded-authority confirmation & Live pilot, decision-impact audit, tokens, failures, and confirmation lock \\
2026-08-24: source-outcome falsification & Three frozen routes, 100 target seeds per route, 19 matched outcome permutations, and randomization-test outputs \\
2026-08-24: hidden-target non-interference & Exact shadow-label reconstruction for 170 archived decision states, route-level summaries, and verified prompt hashes \\
2026-08-24: LLM probability calibration & All 170 archived decisions, fixed-bin reliability data, same-candidate GP comparison, and route-bootstrap uncertainty \\
2026-08-24: FLIP2 external task family & Frozen official wild-type split, 100-seed GP/RGPE/ICM/skill-prior comparison, protocol lock, one live Opus trace, and negative-transfer audit \\
2026-08-24: FLIP2 source-only family gate & Two 50-seed pseudo-target calibration routes, frozen eligibility calculation, outcome-blind P06241 deployment replay, protocol lock, and real-data gate matrix \\
2026-08-25: FLIP2 Rhodopsin prospective gate & Commit-before-download preregistration, two 50-seed source-only calibration routes, 100 untouched official-test seeds, per-seed audits, protocol locks, and real-data confidence-interval figure \\
2026-08-25: FLIP2 IRED prospective positive transfer & Outcome-blind Opus hypothesis, commit-before-download preregistration, 100 calibration seeds, 100 untouched official-test seeds, five matched policies, 6,000 test audit rows, bootstrap statistics, and real-data trajectory figure \\
2026-08-25: IRED source-outcome falsification & One true binding, 99 deterministic outcome-permuted bindings, 100 matched target seeds, 10,000 skill trajectories, 100 target-GP trajectories, assignment-level randomization test, and source-data figure \\
2026-08-25: FLIP2 TrpB prospective replication & Commit-before-download Opus hypothesis, outcome-blind sequence-hash pool, 100 calibration and 100 deployment seeds, strong matched baselines, 99 preregistered source-outcome permutations, 10,100 mechanism trajectories, and source-data figure \\
\bottomrule
\end{tabularx}
\end{table}

\section{Prompt and trace transparency}
\label{sec:prompts}

\subsection{Offline semantic-skill compiler}

\begin{lstlisting}[style=carecode]
INPUT
  source task description and permitted source evidence
  public target fields, categorical values, and counts
  permitted rule operators and acquisition policies

OUTPUT
  skill_id and falsifiable rationale
  executable semantic feature rules
  optional coefficient direction and strength
  exploration/acquisition schedule
  explicit assumptions and failure conditions

VALIDATION
  reject unknown fields, values, operators, or malformed schema
  calibrate on development seeds only
  freeze before held-out replay
\end{lstlisting}

\subsection{Prospective outcome-blind IRED hypothesis}

\begin{lstlisting}[style=carecode]
MODEL
  anthropic/claude-opus-4-8, temperature 0

VISIBLE INPUT
  public IRED task description and official partition counts
  source: variants with 0--2 mutations
  target schema: variants with >2 mutations
  endpoint: best-so-far AUC under 3 + 12 target reveals
  menu: additive-mutation skill or abstain

HIDDEN INPUT
  raw dataset and all train, validation, and test activities

OUTPUT
  recommended_skill: source_additive_mutation_prior
  hypothesis: low-order mutation effects are approximately additive
  mechanism: shrunk token effects + class fallback + decaying authority
  failure conditions: epistasis, unseen positions, distribution shift,
                      excessive early source authority

AUDIT
  request_sha256: 0a42117a5e54c048...f311cccb5e23f78a47f833
  usage: 636 prompt + 775 completion = 1,411 tokens
  generated and committed before dataset download
\end{lstlisting}

The stored record contains the complete system and user messages, unmodified
model response, parsed JSON, request hash, token usage, and the explicit flag
\code{target\_outcomes\_available=false}. The generated hypothesis is useful
only because it is converted into an executable skill and evaluated against
matched target-only and transfer baselines; the text alone is not counted as
evidence.

\subsection{Online proposer--critic record}

\begin{lstlisting}[style=carecode]
{
  "round": 1,
  "hypothesis": "falsifiable source-to-target relation",
  "continue_source_transfer": true,
  "candidate_menu": ["target_gp", "source_prior", "consensus", "coverage"],
  "proposer_choice": "candidate_id",
  "critic_choice": "candidate_id",
  "selected_candidate": "candidate_id",
  "selected_gp_rank": 4,
  "revealed_target_value": 86.875,
  "fallback_reason": null,
  "token_usage": {"input": 0, "output": 0}
}
\end{lstlisting}

The actual traces store model-specific token values and the complete prompts;
the zeros above are schema placeholders rather than experimental measurements.
An explanation is counted as decision-relevant only when it changes the
executed candidate, transfer state, or handoff policy. This rule prevents a
verbose trace from being mistaken for causal contribution.

\section{Limitations and claim boundary}
\label{sec:limitations}

The evaluation is finite-pool replay on public measured datasets. It preserves
experimental budgets and hides unrevealed outcomes, but it cannot model assay
failure, batch effects, synthesis feasibility, or human intervention in a new
laboratory campaign. Most online routes currently have one successful model
trajectory. Several routes share task families, and controller policies were
developed with knowledge of earlier trajectories. The live bounded-authority
pilot had zero decision impact. The completed 100-seed source-outcome
falsification did not reject matched outcome-assignment nulls on any of three
frozen routes. The frozen external FLIP2 task family produced significant
negative transfer for all three numerical transfer baselines; the LLM
abstained online but its initial design was worse than the fixed comparator.
On a second FLIP2 family, a commit-before-download source-only gate
prospectively rejected all three transfer baselines and avoided the 1.548 AUC
loss of the ungated choice. A third FLIP2 family supplied one prospective
positive result: the outcome-blind LLM additive skill improved official-test
AUC by 0.389 [0.274, 0.503] against target-only GP-UCB. The predeclared fixed
skill prior improved by 0.864 and the LLM skill did not improve terminal best
with a confidence interval excluding zero. A later post-hoc IRED assignment
audit found that the true source binding exceeded the mean of 99 false bindings
by 0.340 AUC [0.259, 0.422], but six null assignments were at least as large
($p=0.07$). This is one external finite-pool efficacy result with suggestive,
not conclusive, mechanism attribution; it is not a general LLM advantage.

The independently frozen TrpB protocol strengthens that boundary. Its
LLM-selected additive skill improved official-test AUC by 0.949 and its
predeclared source-assignment test rejected the 99-permutation null at
$p=0.01$. This shows that a model-generated executable skill can carry measured
source information into a new enzyme family under a fixed budget. The same
experiment also falsifies a stronger complete-system claim: the confidence-
bound router abstained because its single-mutant calibration route had no exact-
token coverage, so the deployed policy remained target-only. No result should
therefore be described as universal
cross-domain transfer or prospective wet-lab discovery.

The next decisive evidence is: (i) completion of the frozen repeated-call
protocol, (ii) replication in a more distant domain with stronger fixed and
learned baselines, (iii) a coverage-aware router frozen on a new family, and
(iv) at least one paired
prospective wet-lab campaign. Until then, the principal contribution is an auditable
method for compiling, testing, and rejecting scientific transfer skills.

\end{appendices}
\end{document}
